\section{Tasks in AI Software Engineering} \label{sec:main-tasks-milestones}
We first provide a taxonomy of tasks in AI software engineering. To provide a structured way to consider each task, we define three measures that apply across them: scope, logical complexity, and level of human intervention. To achieve an AI software engineer, we strive for AI to be capable across the board for all three measures.

\textbf{Scope Measure}: We define three levels of scope, the extent of changes that the AI makes to the codebase. \textit{Function-level} scope refers to single, self-contained functions such as in HumanEval \citep{chen2021evaluating} and MBPP \citep{austin2021program}. \textit{Self-contained unit} scope goes beyond singular functions and to larger chunks of code such as entire files and classes. \textit{Project-level} scope refers to larger codebases such as entire repositories, such as in Commit0 \citep{zhao2024commit0} and SWE-Bench \citep{jimenez2024swebench}.

% The AI for code community generally began with function-level benchmarks, where the goal was to write or complete a single self-contained function. In early function-level benchmarks, the specification generally contained a natural language description of the function and test cases. Examples of this are HumanEval \citep{chen2021evaluating} and MBPP \citep{austin2021program}. Later, to mimic more real-world scenarios, function-level benchmarks began to require usage of more specialized APIs and packages such as \texttt{numpy, pandas}, and \texttt{scipy}. Some of these benchmarks include ODEX \citep{wang2022execution}, Jigsaw \citep{jain2022jigsaw}, DS-1000 \citep{lai2023ds}, and JuICe \citep{agashe2019juice}, and BigCodeBench \citep{zhuo2024bigcodebench}.
% \item \textbf{Self-contained unit}: Beyond singular functions, people have studied coding capabilities in self-contained units like files, classes 
% \item \textbf{Repository-level}: With increasing capabilities of AI systems, repository-level 
% Commit0 \citep{zhao2024commit0}, SWE-Bench \citep{jimenez2024swebench}
% \end{itemize}

\textbf{Logical Complexity Measure}: Tasks require a wide range of reasoning abilities when it comes to devising algorithms to solve them. \textit{Low logical complexity} tasks require little to no reasoning, such as writing CRUD (create, read, update, delete) applications or using APIs. \textit{Medium logical complexity} tasks include most LeetCode problems, finding inputs to fuzz simple programs, and reasoning about execution behavior of multithreaded programs. \textit{High logical complexity} tasks require meticulous and challenging levels of algorithmic and logical reasoning, either because the algorithm is complex or because the problem requires clever thinking and insights. This includes difficult competition programming problems, writing large thread-safe concurrent programs, cracking cryptographic ciphers, and solving SMT-like problems. Many popular coding benchmarks are function-level, medium-high logical complexity, such as APPS \citep{hendrycksapps2021}, CodeContests \citep{li2022competition}, and LiveCodeBench \citep{jain2024livecodebenchholisticcontaminationfree}. 

\textbf{Level of Human Intervention Measure}: AI coding is a collaborative task. We follow the autonomy taxonomy outlined in \citet{morris2023levels} to define three levels of human intervention. \textit{Low autonomy} is when the human has full control over the task and uses AI to automate simple sub-tasks. \textit{Medium autonomy} is when the human and AI contribute a similar amount, with interactive coordination of goals and tasks. Here, the human and AI might both suggest refactorings and optimizations during the development cycle. \textit{High autonomy} is when AI drives the interaction and tasks, identifying required changes and the changing demands of the user. The AI would defer to the human only when needed or for a final check, write the code and tests autonomously. Next, we turn to the set of tasks that are reflective of the tasks and capabilities of a human software engineer. 
% \ms{maybe add 1-2 examples for each? code completion is "low" ig}

% \todo{Talk about a few benchmarks or tools, and fit them on these three measures}.


% \todo{Algorithmic complexity vs Repository Scope (ex: CRUD app, SMT, Cipher). OOD stuff in high algorithmic complexity.}

% \todo{Call these Measures (scale, complexity, lvl of human interv)}

% \textbf{Scope}: Tasks in the field of AI for code vary by scope. Here, we categorize scope into several groups:
% \begin{itemize}

%     \item \textbf{Algorithmic and competition programming}: With the increased attention to the reasoning abilities of large language models, competition programming has emerged as an attractive domain to evaluate coding capabilities. Examples of competitive programming benchmarks include APPS \citep{hendrycksapps2021}, CodeContests \citep{li2022competition}, USACOBench \citep{shi2024can}, Code4Bench \citep{majd2019code4bench}, \todo{more}, and LiveCodeBench \citep{jain2024livecodebenchholisticcontaminationfree}. This domain has been attractive because problems are often crisp, concise, and have no ambiguity. Solving these problems requires sharp algorithmic reasoning capabilities, often necessitating insights that may not be obvious upon first reading the problem. Solutions are generally written in a single file comprising of at most a few functions and rarely requiring external packages beyond the standard library. 
%     \item \textbf{Self-contained unit}: Beyond singular functions, people have studied coding capabilities in self-contained units like files, classes 
%     \item \textbf{Repository-level}: With increasing capabilities of AI systems, repository-level 
%     Commit0 \citep{zhao2024commit0}, SWE-Bench \citep{jimenez2024swebench}
%     \item OOD: This could include niche repositories, domains, low-resource languages. CRUXEval, Test Generation, Vulnerability Detection, Code Optimization
% \end{itemize}


% \todo{diff domains have diff requirements for the 3 measures}

% \textbf{Domains}: Web, Games, Security ??
% Web: low reasoning complexity
% SMT: high reasoning complexity

% \todo{Strengths and weaknesses of a few interesting benchmarks}

\subsection{Code Generation}

Code generation is the task of generating code from a specification. In \textit{code completion}, the specification takes the form of a preexisting code snippet, and the goal is to complete the snippet. There are two popular paradigms for code completion: \textit{tab completion}, where the user can press the tab key to complete a block of code (e.g. GitHub Copilot), and \textit{natural language to code}, where the specification is a natural language description with requirements such as the task description or input-output examples. Recently, AI-driven IDEs have blurred the lines between the two paradigms. With the ultimate goal of decreasing the burden of human programmers, they aim to automatically infer the user's intent from the code context and user behavior. However, when intent is vague, they allow users to specify desired functionality via chat interfaces.


\subsection{Code Transformation}

\textbf{Code Refactoring}: In code refactoring, the goal is to take a working implementation of a piece of software and rewrite parts of it while maintaining correctness. One challenge with this task is that success extends beyond functional correctness or metrics. Because it can often be unclear what level of abstraction refactorings should be done at, completing a refactoring at a high autonomy level is also difficult. These challenges are further compounded by the need to understand implicit trade-offs customized to specific codebases, respect conventions, and reason about the long-term maintenance implications of structural changes. 

\textbf{Code Migration and Translation}: An incredibly resource-intensive task is migrating large amounts of code while preserving all the original functionality. Such high-value migrations present opportunities for AI-assisted automation to reduce cost and manual effort. Code migration often involves changes across many files and systems with complex transformations. A special case of code migration is code translation (transpilation): rewriting code from a source language to a target language. In industry, this task can be motivated by several reasons such as security and scalability concerns in legacy languages or avoiding technical debt. Due to the safety-critical and cross-system nature of many migrations, this task often requires substantial human oversight in practice and cannot be done fully autonomously. 

\textbf{Code Optimization}: Transforming programs to improve performance characteristics while maintaining functional correctness is a critical software task. Optimizing real-world systems poses significant challenges, as performance bottlenecks must be identified and new algorithms to mitigate them must be proposed. Code optimization often has a large search and solution space with competing objectives like speed, memory efficiency, and readability, for example when optimizing kernel code at the PTX level for GPU-based AI model optimization~\citep{deepep2025, ouyang2024kernelbench}. 

\subsection{Software Testing and Program Analysis}

\textbf{Software Testing}: Software testing is a practical approach to prevent bugs, both during development and production. There are several popular approaches to software testing, some short-term and others longer-term. For example, \textit{Unit testing} refers to using input-output style tests that exercise the functionality of a piece of code and \textit{property-based testing} is based on formal specifications and relies on specifying test cases that ensure that known properties of the code hold. The goal of software testing is to design tests that can surface bugs reliably. This is evaluated through metrics such as code coverage--how much of the source code is executed when the test suite is run. While practical, software testing faces challenges such as the scalability limits of traditional tools and the difficulty of manually designing tests with good coverage. 

\textbf{Program Analysis}: While testing catches bugs, the most challenging software issues are security vulnerabilities and zero-day exploits, from memory corruption to privilege escalation. This requires a deep program understanding, that testing/fuzzing often misses. For instance, a \textit{zero-day} is a vulnerability unknown to the software developers that is found by an attacker, and there is no patch available from the vendor. In such cases, the only practical approach is offensive security research, manual source code audits, and root cause analysis of prior vulnerabilities to harden codebases.

\subsection{Software Maintenance}

\textbf{Code Documentation and Summarization}: To ensure maintainability, readability, and ease of collaboration, code must be well documented. Good documentation needs to be written cleanly and crisply, describing what the function does and how the function works. It must also anticipate and address any misunderstandings that a programmer might have, such as potential side effects or special cases. Humans often see documentation as a chore and neglect it, leading to code and documentation frequently being out of sync. This has led to the concept of ``self-documenting code'', code that clearly conveys its purpose. 

\textbf{Pull Request (PR) Review}: Reviewing pull requests is an integral aspect of the software development cycle. While the most essential requirement for PRs is that a new feature is implemented correctly, other important considerations include checking whether the repository's style conventions are satisfied, ensuring that the PR does not introduce any new bugs, verifying that program invariants and guarantees still hold, and inspecting whether tests are robust. 

\textbf{Code Understanding, Navigation, and Question Answering}: When encountering a codebase for the first time, developers often find it challenging to understand and develop a good mental model of the code. This can be due to many reasons: too many wrapper functions, excessive error-handling boilerplate, deep call stacks, or poor code cleanliness. One important challenge in code understanding is \textit{code navigation}: finding where relevant functionality is implemented. Doing this well requires understanding the high-level layout of where every functionality lies in the codebase and the low-level understanding of which helper functions are used to implement each functionality. 

Another challenge is \textit{code question answering}: answering complex questions about a codebase, which requires sophisticated code understanding and reasoning abilities. Models should not hallucinate or give incorrect information that skews a developer's mental model of the code. Beyond other tasks mentioned in this section, developers might commonly ask questions related to data flow (when and where data structures get mutated), code functionality (whether there are any side effects), performance characteristics (determining the runtime and memory complexity of a function), or error handling (whether certain corner cases are handled). 


\subsection{Scaffolding and Meta-Code}

For a software system to work, the core logic must be written well, but that is not enough. Many infrastructural aspects must be in place to support the software. We group these into two main categories: \textit{scaffolding} and \textit{meta-code}. We define \textit{scaffolding} as a task outside of the code that must be done to get the software running properly. Examples of scaffolding include setting up Google authentication, subscribing to APIs, managing file storage, and generating API tokens. In contrast, we define \textit{meta-code} to be code that is important to make the system work but does not actually participate in the execution of its main logic. Examples of meta-code include test harnesses, files, CI/CD code, Makefiles, Dockerfiles, sandbox databases, and preprocessors. Scaffolding and meta-code often are small in scope and have low logical complexity but can require a lot of domain-specific knowledge about the application, requiring human intervention.

\subsection{Formal Verification}

The task of formal verification involves generating checkable, mechanized proofs that can guarantee that a piece of code works as intended. Formal verification of software is necessary in mission-critical applications such as aircraft software, where it is crucial that code is correct with absolute certainty. Over the years, there have been countless programming languages designed specifically for formal verification. Some of the popular ones include TLA \citep{lamport1994introduction}, Coq \citep{Coq-refman}, Lean \citep{de2015lean}, Dafny \citep{leino2010dafny}, Isabelle \citep{nipkow2002isabelle}, and Verus \citep{lattuada2024verus}. While formal verification tools have begun to see adoption in industry, they has not yet become mainstream because of these challenges. Code LLMs could greatly ease this burden and make it more feasible to verify code at larger scales, especially verifying properties requiring lower logical complexity.